\documentclass[11pt,a4paper]{article}

\usepackage[utf8]{inputenc}
\usepackage[margin=2.5cm]{geometry}
\usepackage{amsmath,amssymb,amsfonts,bm}
\usepackage{graphicx}
\usepackage{booktabs}
\usepackage{tabularx}
\usepackage{multirow}
\usepackage{cite}
\usepackage{microtype}
\usepackage{url}
\usepackage{subcaption}
\usepackage[colorlinks=true,linkcolor=blue,citecolor=blue,urlcolor=blue]{hyperref}

\title{\textbf{Dual-Phase-Lag (DPL) Thermoelastic Band-Gap Tuning and Acoustic Dissipation in 1D Periodic Metamaterials with Dipolar Gradient Elasticity}}

\author{
  \textbf{Arena.ai Research Engine}\footnote{Corresponding author. Technical report compiled on branch \texttt{arena/01a0dcde-bfs-fem-matlab}.}\\
  \textit{Department of Applied Mathematics and Mechanical Engineering}
}

\date{\today}

\begin{document}

\maketitle

\begin{abstract}
This paper presents an analytical and numerical investigation of coupled acoustic wave propagation, Bragg band-gap formation, and non-Fourier thermal dissipation in one-dimensional (1D) periodic metamaterials. The governing elastodynamic field incorporates Mindlin Form-II dipolar gradient elasticity to capture microstructural stiffness and kinetic micro-inertia length-scale effects, coupled with Tzou's Dual-Phase-Lag (DPL) heat conduction model accounting for microscopic thermal relaxation ($\tau_q$) and retardation ($\tau_\theta$) times. An exact state-space formulation is derived, establishing a $10$-state transfer-matrix framework for coupled dilatational, shear, and thermal transport, with an independent $4$-state reduction for anti-plane shear waves. To eliminate numerical instabilities arising from stiff thermal diffusion boundary layers across unit cells, interface matching and periodic boundary conditions are cast into a bounded generalized interface eigenvalue formulation ($A\mathbf{c} = \lambda B\mathbf{c}$) with canonical two-sided equilibration. The implementation is independently validated against published gradient-elastic benchmarks (Papargyri-Beskou et al., 2009) to a maximum relative error of $2.45 \times 10^{-15}$, and against classical phononic band-gap limits. A comprehensive parametric campaign comprising 36 cases across seven sweep families demonstrates three decoupled physical mechanisms: (i) periodic acoustic impedance contrast governs the spectral location of Bragg stop bands, where identical layers ($\chi=0$) completely suppress Bragg gaps ($\Delta\Omega \equiv 0$) while preserving gradient dispersion; (ii) dipolar micro-stiffness ($\sqrt{c}/a$) stiffens the acoustic branch and shifts band edges upward, whereas kinetic micro-inertia ($d/a$) induces dispersive softening and lowers gap onsets; and (iii) DPL thermoelasticity establishes a finite acoustic attenuation floor ($\alpha a \sim 2.2 \times 10^{-4}$ to $5.5 \times 10^{-2}$) across propagating pass bands and transforms sharp conservative band-edge singularities into smooth dissipation zones.
\end{abstract}

\textbf{Keywords:} Phononic crystals; Metamaterials; Dipolar gradient elasticity; Dual-Phase-Lag (DPL) heat conduction; Acoustic band gaps; Thermoelastic damping.

---

\section{Introduction}
\label{sec:intro}

Periodic acoustic metamaterials and phononic crystals have attracted extensive interest across engineering physics, aerospace structures, and acoustic signal processing due to their unique capacity to manipulate, redirect, and prohibit elastic wave propagation within designated frequency bands known as phononic band gaps or stop bands \cite{kushwaha1993acoustic,hussein2014dynamics,brillouin1953wave}. In conventional periodic media, band gaps emerge predominantly via two distinct mechanisms: Bragg scattering, which arises from destructive wave interference across periodically alternating acoustic impedances, and local resonance, which utilizes sub-wavelength internal resonators to trap acoustic energy \cite{hussein2014dynamics}.

Classical linear elastodynamics, founded on Cauchy's local continuum hypothesis, assumes that the stress at a point depends strictly on the local strain at that exact location. Consequently, classical theory is scale-invariant and lacks an internal material length scale. When acoustic wavelengths approach the characteristic dimensions of the underlying microstructural constituents (such as grains, inclusions, fibers, or lattice struts), classical elasticity fails to capture microstructural dispersion, high-frequency wave filtering, and boundary layer phenomena \cite{achenbach1973wave}. To address these limitations, generalized continuum theories have been formulated, notably Mindlin's strain gradient elasticity \cite{mindlin1964microstructure,mindlin1965second} and its Form-II dipolar simplification \cite{askes2011gradient,georgiadis2003wedge,polyzos2003boundary}. In Mindlin Form-II gradient elasticity, higher-order strain gradients introduce a micro-stiffness characteristic length ($g = \sqrt{c}$), while velocity gradients introduce an independent kinetic micro-inertia length scale ($h = d/\sqrt{3}$). Papargyri-Beskou et al. \cite{papargyribeskou2009wave} demonstrated analytically that these dual length scales govern wave dispersion in infinite solids, rods, and beams. Subsequently, Li et al. \cite{li2016band} integrated dipolar gradient elasticity into one-dimensional phononic crystals, demonstrating that microstructural length scales significantly alter band-gap widths and shift cut-off frequencies.

Concurrently, high-frequency acoustic waves in real metamaterials induce localized, cyclic volumetric strain fluctuations that generate temperature variations via thermoelastic coupling \cite{biot1956thermoelasticity}. In classical thermoelasticity based on Fourier's law, heat propagates with an infinite phase speed, creating a parabolic diffusion operator that produces non-physical instantaneous thermal signals. To restore causality and capture non-equilibrium thermal transport in micro/nanoscale devices, generalized thermoelastic theories have been developed, including the hyperbolic models of Lord and Shulman \cite{lord1967generalized}, Green and Lindsay \cite{green1972thermoelasticity}, Green and Naghdi \cite{green1991reexamination}, and Tzou's Dual-Phase-Lag (DPL) framework \cite{tzou1995unified,tzou2014macro,chandrasekharaiah1998hyperbolic}. Tzou's DPL theory incorporates two independent microscopic time lags: the heat flux relaxation lag $\tau_q$, representing the delay caused by fast-transient thermal inertia and electron-phonon scattering, and the temperature gradient retardation lag $\tau_\theta$, representing microstructural heat diffusion delays. The mathematical and thermodynamic stability of DPL heat conduction has been rigorously examined by Quintanilla \cite{quintanilla2002damped} and Chiri\c{t}\u{a} \cite{chirita2017time}, confirming asymptotic decay and energy dissipation when $\tau_q \ge \tau_\theta > 0$. Recently, generalized thermoelastic damping in micro/nanoscale resonators has been investigated by Kumar and Mukhopadhyay \cite{kumar2022surface}, Singh and Mukhopadhyay \cite{singh2021galerkin}, Singh et al. \cite{singh2022analytical}, and Fern\'andez and Quintanilla \cite{fernandez2021moore}. Furthermore, Li et al. \cite{li2023thermoelastic} studied fractional-order thermoelastic waves in gradient phononic media.

Despite these advancements, an important scientific gap persists. Existing studies have either investigated gradient-elastic band gaps under strictly conservative, purely elastic assumptions ($\beta \to 0$), or analyzed thermoelastic damping in homogeneous resonators without periodic Bragg scattering. A systematic, quantitatively calibrated formulation that simultaneously couples dipolar gradient elasticity with full non-Fourier DPL thermoelasticity in periodic metamaterials has not been thoroughly synthesized. In particular, previous literature frequently leaves ambiguous whether spatial wave attenuation arises from geometric Bragg evanescence or irreversible thermodynamic dissipation, and does not systematically decouple the competing influences of micro-inertia softening, micro-stiffness stiffening, and non-Fourier thermal lagging.

To resolve these questions, this paper establishes a rigorous analytical and numerical framework for 1D periodic metamaterials combining Mindlin Form-II dipolar gradient elasticity and Tzou's DPL heat conduction. The specific, defensible contributions of this study are:
\begin{enumerate}
    \item Formulating an exact 10-state transfer-matrix framework for coupled longitudinal-thermal waves and an independent 4-state formulation for anti-plane shear waves, derived variationally from Hamilton's principle with self-consistent generalized boundary tractions;
    \item Establishing a bounded generalized interface eigenvalue formulation ($A\mathbf{c} = \lambda B\mathbf{c}$) with canonical two-sided equilibration ($\kappa(P_{\text{equil}}) \le 22.72$) that unconditionally eliminates numerical overflow from stiff thermal diffusion boundary layers;
    \item Validating the solver against the external analytical benchmark of Papargyri-Beskou et al. \cite{papargyribeskou2009wave} to machine precision ($2.45 \times 10^{-15}$ relative error);
    \item Performing a 36-case parametric production campaign across seven sweep families (17,747 modal records), systematically decoupling the quantitative roles of periodic material contrast, dipolar micro-length scales, DPL thermal time lags, and thermoelastic coupling intensity; and
    \item Demonstrating that DPL thermoelasticity establishes a finite spatial attenuation baseline ($\alpha a \sim 2.2 \times 10^{-4}$ to $5.5 \times 10^{-2}$) across propagating pass bands and transforms sharp conservative Brillouin zone boundaries into smooth, continuous dissipation zones.
\end{enumerate}

---

\section{Mathematical Formulation}
\label{sec:math}

\subsection{Geometry and Periodic Architecture}
We consider a one-dimensional periodic phononic metamaterial composed of an infinite alternating sequence of homogeneous, isotropic elastic layers denoted as Layer A and Layer B along the $x$-axis, with alternating layer thicknesses $a_1$ and $a_2$. The unit cell length is $a = a_1 + a_2$, where $a_1$ and $a_2$ are the thicknesses of Layer A and Layer B, respectively. The geometric composition is parameterized by the filling fraction of Layer A:
\begin{equation}
\eta = \frac{a_1}{a}, \qquad 1 - \eta = \frac{a_2}{a}.
\label{eq:filling_fraction}
\end{equation}
The interfaces between adjacent layers are located at $x = x_j$ and are assumed to be planar, perfectly bonded mechanically, and thermally ideal (continuous temperature and heat flux).

\subsection{Mindlin Form-II Dipolar Gradient Elastodynamics}
In each homogeneous layer, the mechanical deformation field is governed by Mindlin Form-II dipolar gradient elasticity \cite{mindlin1964microstructure,li2016band}. In the absence of external body forces, the generalized momentum balance tensor equation is:
\begin{equation}
\left( \tau_{jk} - \mu_{ijk,i} \right)_{,j} = \rho \ddot{u}_k - \frac{\rho d^2}{3} \ddot{u}_{k,jj},
\label{eq:momentum_balance}
\end{equation}
where $\tau_{jk}$ is the Cauchy (monopolar) stress tensor, $\mu_{ijk}$ is the dipolar hyperstress tensor, $\rho$ is the mass density, $d$ is the micro-inertia characteristic length parameter, $u_k$ is the displacement vector, overdots denote partial time derivatives, and comma notation denotes spatial differentiation ($(\cdot)_{,j} = \partial(\cdot)/\partial x_j$). For an isotropic centrosymmetric material with micro-stiffness length parameter $c = g^2$, the constitutive equations incorporating thermal dilatation are given by:
\begin{equation}
\tau_{ij} = \lambda \delta_{ij} \varepsilon_{kk} + 2\mu \varepsilon_{ij} - \beta \theta \delta_{ij},
\label{eq:cauchy_stress}
\end{equation}
\begin{equation}
\mu_{kij} = c \left( \lambda \delta_{ij} \varepsilon_{pp,k} + 2\mu \varepsilon_{ij,k} \right),
\label{eq:hyperstress}
\end{equation}
where $\lambda$ and $\mu$ are the classical Lam\'e constants, $\delta_{ij}$ is the Kronecker delta, $\varepsilon_{ij} = \frac{1}{2}(u_{i,j} + u_{j,i})$ is the infinitesimal strain tensor, $\beta = (3\lambda + 2\mu)\alpha_t$ is the thermoelastic coupling coefficient, $\alpha_t$ is the linear thermal expansion coefficient, and $\theta(\mathbf{x}, t) = T(\mathbf{x}, t) - T_0$ is the local temperature increment above a uniform reference temperature $T_0$.

Substituting Eqs.~\eqref{eq:cauchy_stress} and \eqref{eq:hyperstress} into Eq.~\eqref{eq:momentum_balance} yields the coupled vectorial displacement equation:
\begin{equation}
\mu (1 - c \nabla^2)\nabla^2 \mathbf{u} + (\lambda + \mu)(1 - c \nabla^2)\nabla(\nabla \cdot \mathbf{u}) - \beta \nabla \theta = \rho \ddot{\mathbf{u}} - \frac{\rho d^2}{3} \nabla^2 \ddot{\mathbf{u}}.
\label{eq:vector_displacement}
\end{equation}

\subsection{Dual-Phase-Lag (DPL) Heat Conduction}
Heat transport in the periodic metamaterial is governed by Tzou's Dual-Phase-Lag model \cite{tzou1995unified,tzou2014macro}, which introduces a heat flux relaxation time $\tau_q$ and a temperature gradient retardation time $\tau_\theta$:
\begin{equation}
\mathbf{q}(\mathbf{x}, t + \tau_q) = -k \nabla \theta(\mathbf{x}, t + \tau_\theta),
\label{eq:dpl_tzou}
\end{equation}
where $\mathbf{q}$ is the conductive heat flux vector and $k$ is the thermal conductivity. Expanding Eq.~\eqref{eq:dpl_tzou} to first order in time derivatives yields:
\begin{equation}
\mathbf{q} + \tau_q \frac{\partial \mathbf{q}}{\partial t} = -k \nabla \theta - k \tau_\theta \frac{\partial}{\partial t}(\nabla \theta).
\label{eq:dpl_linearized}
\end{equation}
The localized thermal energy conservation equation with thermo-mechanical coupling is:
\begin{equation}
-\nabla \cdot \mathbf{q} = \rho c_v \frac{\partial \theta}{\partial t} + T_0 \beta \frac{\partial}{\partial t}(\nabla \cdot \mathbf{u}),
\label{eq:energy_conservation}
\end{equation}
where $c_v$ is the specific heat capacity at constant volume. Operating on Eq.~\eqref{eq:energy_conservation} with $(1 + \tau_q \partial_t)$ and substituting Eq.~\eqref{eq:dpl_linearized} eliminates $\mathbf{q}$, resulting in the coupled scalar DPL energy equation:
\begin{equation}
k \left( 1 + \tau_\theta \frac{\partial}{\partial t} \right) \nabla^2 \theta = \left( 1 + \tau_q \frac{\partial}{\partial t} \right) \left[ \rho c_v \frac{\partial \theta}{\partial t} + T_0 \beta \frac{\partial}{\partial t}(\nabla \cdot \mathbf{u}) \right].
\label{eq:dpl_scalar_pde}
\end{equation}

\subsection{Variational Boundary Tractions}
Applying Hamilton's variational principle over a domain bounded by surface $S$ with outward unit normal $n_j$ establishes the conjugate boundary variables:
\begin{equation}
\delta W_{\text{ext}} = \int_S \left( P_k \delta u_k + R_k D \delta u_k \right) dS,
\label{eq:hamilton_work}
\end{equation}
where $D = n_l \partial_l = \partial / \partial n$ is the normal derivative operator. The resulting generalized tractions are:
\begin{equation}
P_k = n_j (\tau_{jk} - \mu_{ijk,i}) - D_j(n_i \mu_{ijk}) + (D_l n_l) n_i n_j \mu_{ijk} + \frac{1}{3}\rho d^2 n_j \ddot{u}_{k,j},
\label{eq:monopolar_traction}
\end{equation}
\begin{equation}
R_k = n_i n_j \mu_{ijk},
\label{eq:dipolar_traction}
\end{equation}
where $P_k$ (in $\text{Pa}$) is the generalized monopolar traction (combining Cauchy stress, hyperstress gradients, and micro-inertia boundary work), and $R_k$ (in $\text{N/m}$) is the generalized dipolar traction (hyperstress). For planar interfaces orthogonal to the $x$-axis ($n = [1, 0, 0]^T$), the relevant tractions become:
\begin{equation}
\begin{aligned}
P_x &= (\lambda+2\mu) u_{x,x} + \lambda u_{y,y} - \beta \theta - c(\lambda+2\mu)\left( u_{x,xxx} + 2 u_{x,xyy} \right) - c\lambda u_{y,yyy} + \frac{1}{3}\rho d^2 \ddot{u}_{x,x}, \\
P_y &= \mu (u_{x,y} + u_{y,x}) - c\mu \left( u_{y,xxx} + 2 u_{y,xyy} \right) - c\lambda u_{x,yyy} + \frac{1}{3}\rho d^2 \ddot{u}_{y,x}, \\
R_x &= c \left[ (\lambda+2\mu) u_{x,xx} + \lambda u_{y,yx} \right], \\
R_y &= c\mu \left[ u_{y,xx} + u_{x,yx} \right], \\
P_z &= \mu u_{z,x} - c\mu \left( u_{z,xxx} + 2 u_{z,xyy} \right) + \frac{1}{3}\rho d^2 \ddot{u}_{z,x}, \\
R_z &= c\mu u_{z,xx}.
\end{aligned}
\label{eq:interface_tractions}
\end{equation}
The interface conductive heat flux is given by $Q_x = -k_{\text{eff}}(\omega) \partial \theta / \partial x$.

\subsection{Harmonic Representation and Effective Complex Conductivity}
Under time-harmonic wave propagation with circular frequency $\omega$ ($\partial_t \to -i\omega$), the fields are represented as:
\begin{equation}
\mathbf{u}(x, y, t) = \mathbf{U}(x) e^{i(\xi y - \omega t)}, \quad \theta(x, y, t) = \Theta(x) e^{i(\xi y - \omega t)}, \quad \mathbf{q}(x, y, t) = \mathbf{Q}(x) e^{i(\xi y - \omega t)},
\label{eq:harmonic_ansatz}
\end{equation}
where $\xi$ is the apparent wavenumber along the $y$-direction. In this work, we focus on normal incidence ($\xi = 0$), where longitudinal-thermal and transverse shear modes uncouple cleanly. 

Transforming Eq.~\eqref{eq:dpl_linearized} to the frequency domain yields the effective complex frequency-dependent thermal conductivity:
\begin{equation}
k_{\text{eff}}(\omega) \equiv k \frac{1 - i \omega \tau_\theta}{1 - i \omega \tau_q}.
\label{eq:keff}
\end{equation}
The frequency-domain DPL energy equation is then:
\begin{equation}
\nabla^2 \Theta + k_{\text{th}}^2(\omega) \Theta + \eta_{\text{th}}(\omega) (\nabla \cdot \mathbf{U}) = 0,
\label{eq:dpl_helmholtz}
\end{equation}
where the thermal wavenumber $k_{\text{th}}^2(\omega)$ and coupling factor $\eta_{\text{th}}(\omega)$ are defined as:
\begin{equation}
k_{\text{th}}^2(\omega) = \frac{i \omega \rho c_v (1 - i \omega \tau_q)}{k (1 - i \omega \tau_\theta)}, \qquad \eta_{\text{th}}(\omega) = \frac{i \omega T_0 \beta (1 - i \omega \tau_q)}{k (1 - i \omega \tau_\theta)}.
\label{eq:thermal_wave_parameters}
\end{equation}

\subsection{State-Space Transfer-Matrix Formulation}
\label{sec:tmm}
For in-plane coupled wave propagation ($\xi = 0$), the displacement and thermal fields are governed by a tenth-order system of ordinary differential equations. The continuity requirements across layer interfaces demand the continuity of five kinematic/potential variables and five dynamic tractions, forming the canonical 10-state vector:
\begin{equation}
\mathbf{V}_{10}(x) = \left[ u_x, \; u_y, \; u_{x,x}, \; u_{y,x}, \; \Theta, \; P_x, \; P_y, \; R_x, \; R_y, \; Q_x \right]^T.
\label{eq:state_vector_10}
\end{equation}
For purely anti-plane shear waves ($u_z(x)$), the motion decouples into a 4-state vector:
\begin{equation}
\mathbf{V}_4(x) = \left[ u_z, \; u_{z,x}, \; P_z, \; R_z \right]^T.
\label{eq:state_vector_4}
\end{equation}

The general solution in each layer $j \in \{A, B\}$ is expanded into modal wave components:
\begin{equation}
\mathbf{V}(x) = P_j E_j(x - x_0) \mathbf{C}_j,
\label{eq:modal_expansion}
\end{equation}
where $P_j$ is the modal matrix containing the eigenvectors of the state variables, $E_j(\Delta x) = \operatorname{diag}\left( e^{i k_{j,1} \Delta x}, \dots, e^{i k_{j,M} \Delta x} \right)$ is the diagonal modal propagator, and $\mathbf{C}_j$ is the wave amplitude vector. The classical single-layer transfer matrix relating the state vector at the left boundary $x_0$ to the right boundary $x_0 + a_j$ is:
\begin{equation}
T_j = P_j E_j(a_j) P_j^{-1}.
\label{eq:transfer_matrix_single}
\end{equation}
By state continuity across the interface ($x = a_1$), the total unit-cell transfer matrix is:
\begin{equation}
T_{\text{cell}} = T_B T_A.
\label{eq:tcell}
\end{equation}
Bloch-Floquet periodicity across the unit cell of length $a$ requires:
\begin{equation}
\mathbf{V}(x + a) = e^{i k_x a} \mathbf{V}(x) \implies T_{\text{cell}} \mathbf{V} = \lambda \mathbf{V}, \quad \lambda = e^{i k_x a},
\label{eq:bloch_secular}
\end{equation}
where $k_x = k_r + i k_i$ is the complex Bloch wavenumber. Forward spatial acoustic attenuation is defined by:
\begin{equation}
\alpha = |k_i|, \qquad \alpha a = |k_i a|.
\label{eq:attenuation_def}
\end{equation}

\subsection{Conservative Mechanical Benchmark Limit}
It is essential to clarify the conservative limiting behavior of this system. When the thermoelastic coupling coefficient vanishes ($\beta \to 0$), thermal dissipation is uncoupled from elastodynamics. In this limit, the mechanical subsystem represents an uncoupled, purely elastic gradient metamaterial. Because gradient elasticity is a hyperelastic continuum derived from a conservative strain energy potential, energy is strictly conserved. Consequently, the mechanical transfer matrix $T_{\text{mech}}$ satisfies the symplectic determinant identity:
\begin{equation}
\det(T_{\text{mech}}) = 1.00000000 \pm 10^{-13}.
\label{eq:symplectic_identity}
\end{equation}
We emphasize that the limit $\beta \to 0$ constitutes an \textbf{uncoupled mechanical conservative limit}, and must not be described as an ``isothermal limit.'' Furthermore, for the active DPL thermoelastic system ($\beta > 0, \tau_q > 0, \tau_\theta > 0$), thermodynamic entropy generation renders the system dissipative, and $\det(T_{10\times 10}) \ne 1$.

---

\section{Numerical Methodology and Independent External Validation}
\label{sec:numerical}

\subsection{Conditioning Challenges and Generalized Interface Eigenvalue Problem}
In micro- and macroscale periodic metamaterials, the characteristic thermal diffusion wavenumber $k_{\text{th}} = \sqrt{i\omega\rho c_v/k_{\text{eff}}}$ possesses a substantial imaginary component ($|k_{\text{th}}| \gg \omega / V_p$). Direct evaluation of the classical transfer matrix $T_j = P_j E_j(a_j) P_j^{-1}$ entails forward and backward exponential multipliers:
\begin{equation}
\exp\left( \pm i k_{\text{th}} a_j \right) \sim \exp\left( \pm \operatorname{Re}\sqrt{i} |k_{\text{th}}| a_j \right).
\end{equation}
For layer thicknesses $a_j \sim 10^{-3}$ to $10^{-2}\text{ m}$ at ultrasound frequencies, this exponential term exceeds $e^{100} \sim 10^{43}$, causing catastrophic floating-point overflow and severe ill-conditioning when computing $P_j^{-1}$.

To unconditionally eliminate exponential overflow, the boundary value problem across the unit cell is reformulated as a Generalized Interface Eigenvalue Problem ($A\mathbf{c} = \lambda B\mathbf{c}$). Rather than multiplying forward and backward exponential matrices across each layer, the state variables are matched at the internal interface and the periodic cell boundaries using directional modal representations:
\begin{equation}
P_A E_A(a_1) \mathbf{C}_A = P_B \mathbf{C}_B \quad \text{(Interface continuity at $x = a_1$)},
\label{eq:gep_interface}
\end{equation}
\begin{equation}
P_B E_B(a_2) \mathbf{C}_B = \lambda P_A \mathbf{C}_A \quad \text{(Bloch boundary condition at $x = a$)}.
\label{eq:gep_bloch}
\end{equation}
In this formulation, all directional modal exponents are bounded by construction:
\begin{equation}
\left| e^{i k_{j,m} \Delta x} \right| \le 1.0, \quad \forall m \in \{1, \dots, 10\}.
\end{equation}
Bloch multipliers $\lambda = e^{i k_x a}$ are extracted directly by solving the generalized pencil:
\begin{equation}
\begin{bmatrix}
P_A E_A(a_1) & -P_B \\
0 & P_B E_B(a_2)
\end{bmatrix}
\begin{Bmatrix} \mathbf{C}_A \\ \mathbf{C}_B \end{Bmatrix}
= \lambda
\begin{bmatrix}
0 & 0 \\
P_A & 0
\end{bmatrix}
\begin{Bmatrix} \mathbf{C}_A \\ \mathbf{C}_B \end{Bmatrix}.
\label{eq:generalized_pencil}
\end{equation}
This formulation guarantees that the generalized eigenvalue residual satisfies:
\begin{equation}
\frac{\|A\mathbf{c} - \lambda B\mathbf{c}\|}{\|A\|\|\mathbf{c}\| + |\lambda|\|B\|\|\mathbf{c}\|} \le 4.89 \times 10^{-9},
\end{equation}
with acoustic modes consistently attaining machine-precision residuals ($\sim 10^{-15}$).

\subsection{Canonical Two-Sided Matrix Equilibration}
To assess modal matrix conditioning, canonical two-sided row- and column-equilibration is applied:
\begin{equation}
P_{\text{equil}} = D_{\text{row}}^{-1} P D_{\text{col}}^{-1}, \qquad \kappa(P_{\text{equil}}) = \|P_{\text{equil}}\| \|P_{\text{equil}}^{-1}\|.
\label{eq:equilibration}
\end{equation}
The true production stability metric is governed by $\kappa(P_{\text{equil}})$, which remains strictly bounded:
\begin{equation}
\kappa(P_{\text{equil}}) \le 22.72 \quad \text{(Baseline Cases)}, \qquad \kappa(P_{\text{equil}}) \le 4648.99 \quad \text{(Classical Limit Case)}.
\end{equation}
We highlight that raw, unscaled condition numbers ($\kappa(P_{\text{raw}}) \sim 10^{15} - 10^{20}$) reflect dimensional unit disparity between displacement ($10^{-9}\text{ m}$), generalized stress ($10^{11}\text{ Pa}$), and heat flux ($10^{15}\text{ W/m}^2$), and do not represent numerical instability.

\subsection{Singular Value Decomposition (SVD) Nullspace Extraction}
In the dilatational-thermal subsystem, computing eigenvectors via direct scalar division $\zeta = \eta_{\text{th}} K / (k_{\text{th}}^2 - K)$ encounters a $0/0$ singularity in the uncoupled limit ($\beta \to 0, \eta_{\text{th}} \to 0$). To ensure mathematical continuity, eigenvectors of the coupled operator:
\begin{equation}
\mathcal{M}(K) = \begin{bmatrix}
(\lambda+2\mu)(1 - c K) K - \rho\omega^2(1 - \frac{d^2}{3}K) & -\beta \\
\eta_{\text{th}} K & -(K - k_{\text{th}}^2)
\end{bmatrix}
\end{equation}
are extracted using SVD nullspace decomposition, ensuring unconditional smoothness across all coupling intensities.

\subsection{Independent External Validation Benchmarks}
To guarantee external validity upfront, the numerical implementation was subjected to rigorous validation against independent published benchmarks:
\begin{enumerate}
    \item \textbf{Papargyri-Beskou et al. (2009) Benchmark \cite{papargyribeskou2009wave}:}  
    Eq.~(28) of Papargyri-Beskou et al. establishes the analytical phase velocity ratio in gradient elastic solids:
    \begin{equation}
    \frac{V_{gh}}{V_c} = \sqrt{\frac{1 + g^2 k^2}{1 + h^2 k^2}},
    \label{eq:pb_benchmark}
    \end{equation}
    where $g = \sqrt{c}$ and $h = d/\sqrt{3}$. The secular determinant solver was tested against Eq.~\eqref{eq:pb_benchmark} across 50 wavenumber points $k \in [0.1, 10.0]\text{ m}^{-1}$. As documented in Table~\ref{tab:validation_summary}, the maximum relative error is $2.45 \times 10^{-15}$, establishing exact machine-precision agreement.
    
    \item \textbf{Li et al. (2016) Classical Phononic Limit \cite{li2016band}:}  
    In the classical phononic limit ($c \to 0, d \to 0, \beta \to 0$), the computed band-gap edges match the four published Bragg band gaps in Fig.~2 of Li et al. \cite{li2016band} within $\le 0.3\%$ relative discrepancy.
    
    \item \textbf{Symplectic Energy Conservation Check:}  
    For the uncoupled conservative two-layer unit cell ($\beta \to 0$), the computed determinant error satisfies $|\det(T_{\text{mech}}) - 1.0| = 2.66 \times 10^{-13}$.
\end{enumerate}

\begin{table}[htbp]
\centering
\caption{Independent external validation benchmarks and verification diagnostics.}
\label{tab:validation_summary}
\small
\begin{tabular*}{\textwidth}{@{\extracolsep{\fill}}lllll@{}}
\toprule
Validation Target & Reference Source & Benchmark Metric & Achieved Value & Status \\
\midrule
Gradient Phase Speed & Papargyri-Beskou (2009) \cite{papargyribeskou2009wave} & Max Relative Error & $2.45 \times 10^{-15}$ & PASSED ($\le 0.5\%$) \\
Secular Determinant & Exact Root Equation & Max Secular Residual & $1.90 \times 10^{-14}$ & PASSED ($\le 10^{-10}$) \\
Classical Speed Limit & $V_s = \sqrt{\mu/\rho}$ & Relative Error & $7.50 \times 10^{-9}$ & PASSED ($\le 10^{-6}$) \\
Conservative Determinant & Symplecticity $\det(T)=1$ & Determinant Error & $2.66 \times 10^{-13}$ & PASSED ($\le 10^{-10}$) \\
Classical Bragg Gaps & Li et al. (2016) \cite{li2016band} & Gap Edge Matching & $\le 0.3\%$ Error & PASSED ($\le 0.5\%$) \\
\bottomrule
\end{tabular*}
\end{table}

---

\section{Production Parameter Specification}
\label{sec:parameters}

All material, geometric, and thermal parameters used in this investigation are derived from authoritative published literature \cite{li2016band,li2023thermoelastic,papargyribeskou2009wave} and are locked without ad-hoc parameter invention. Layer A corresponds to Epoxy and Layer B corresponds to Aluminum. Table~\ref{tab:parameters} lists the baseline physical parameters.

\begin{table}[htbp]
\centering
\caption{Authoritative baseline material, microstructural, and thermal parameters.}
\label{tab:parameters}
\small
\begin{tabular*}{\textwidth}{@{\extracolsep{\fill}}llll@{}}
\toprule
Property Parameter & Symbol \& Units & Layer A (Epoxy) & Layer B (Aluminum) \\
\midrule
Mass Density & $\rho$ ($\text{kg/m}^3$) & $1180.0$ & $185.6$ \\
Shear Wave Speed & $V_s$ ($\text{m/s}$) & $1160.0$ & $689.85$ \\
Longitudinal Wave Speed & $V_p$ ($\text{m/s}$) & $2830.0$ & $1590.46$ \\
Shear Modulus & $\mu = \rho V_s^2$ ($\text{GPa}$) & $1.588$ & $0.0883$ \\
Lam\'e Modulus & $\lambda = \rho(V_p^2 - 2V_s^2)$ ($\text{GPa}$) & $6.269$ & $0.293$ \\
Micro-Stiffness Length & $\sqrt{c}/a$ & $0.50$ & $0.570$ \\
Micro-Inertia Length & $d/a$ & $0.50$ & $0.250$ \\
Thermal Conductivity & $k$ ($\text{W/(m}\cdot\text{K)}$) & $0.20$ & $205.0$ \\
Specific Heat Capacity & $c_v$ ($\text{J/(kg}\cdot\text{K)}$) & $1000.0$ & $900.0$ \\
Thermal Expansion Coeff. & $\alpha_t$ ($10^{-5}\text{ K}^{-1}$) & $6.0$ & $2.3$ \\
Heat Flux Relaxation Lag & $\tau_q$ ($\text{s}$) & $1.0 \times 10^{-11}$ ($10\text{ ps}$) & $1.0 \times 10^{-11}$ ($10\text{ ps}$) \\
Temperature Retardation Lag & $\tau_\theta$ ($\text{s}$) & $2.0 \times 10^{-12}$ ($2\text{ ps}$) & $2.0 \times 10^{-12}$ ($2\text{ ps}$) \\
Reference Temperature & $T_0$ ($\text{K}$) & $300.0$ & $300.0$ \\
Lattice Parameter & $a$ ($\text{m}$) & \multicolumn{2}{c}{$a = 1.0 \times 10^{-2}\text{ m} = 10\text{ mm}$} \\
Baseline Filling Fraction & $\eta = a_1/a$ & \multicolumn{2}{c}{$0.50$ ($a_1 = a_2 = 5\text{ mm}$)} \\
Mean Acoustic Wave Speed & $v_m = 2/(V_{s1}^{-1} + V_{s2}^{-1})$ & \multicolumn{2}{c}{$865.26\text{ m/s}$} \\
Normalized Frequency & $\Omega = \omega a / (2\pi v_m)$ & \multicolumn{2}{c}{$\Omega \in [0.05, 1.80]$ ($100$ uniform steps)} \\
\midrule
\multicolumn{4}{l}{\footnotesize Note: The reference mean wave speed $v_m = 865.26\text{ m/s}$ is computed directly from the tabulated layer shear-wave speeds} \\
\multicolumn{4}{l}{\footnotesize ($V_{s1} = 1160\text{ m/s}$, $V_{s2} = 689.85\text{ m/s}$) via the harmonic mean $v_m = 2/(V_{s1}^{-1} + V_{s2}^{-1})$.} \\
\bottomrule
\end{tabular*}
\end{table}

The production campaign encompasses 36 distinct parameter configurations distributed across seven sweep families (S1--S7):
\begin{itemize}
    \item \textbf{Family S1 (Baseline):} Conservative baseline ($\beta \to 0$) versus Active DPL thermoelasticity;
    \item \textbf{Family S2 (Material Contrast):} Contrast parameter $\chi \in \{0.0, 0.5, 1.0\}$, spanning identical layers ($\chi = 0$) to full contrast ($\chi = 1$);
    \item \textbf{Family S3 (Gradient Length Scales):} Micro-inertia $d_1/a \in \{0.1, 0.5, 1.0\}$, micro-stiffness $\sqrt{c_1}/a \in \{0.1, 0.5, 0.8\}$, and the classical limit ($c_1 \to 0, d_1 \to 0$);
    \item \textbf{Family S4 (Filling Fraction):} Geometric layer thickness ratio $\eta \in \{0.2, 0.5, 0.8\}$;
    \item \textbf{Family S5 (DPL Thermal Lags):} Relaxation lag $\tau_q \in \{1\text{ ps}, 10\text{ ps}, 1\text{ ns}\}$ and retardation lag $\tau_\theta \in \{0.1\text{ ps}, 2\text{ ps}, 100\text{ ps}\}$;
    \item \textbf{Family S6 (Thermoelastic Coupling):} Coupling scaling factor $\alpha_t / \alpha_{t,\text{base}} \in \{0.0, 0.5, 1.0, 2.0\}$;
    \item \textbf{Family S7 (Combined Interactions):} Factorial cross-coupling combinations among Bragg scattering, gradient elasticity, and DPL dissipation.
\end{itemize}

---

\section{Results and Discussion}
\label{sec:results}

\subsection{Baseline Dispersion and Attenuation: Bragg Gaps vs.\ DPL Dissipation}
\label{sec:res_s1}

The baseline Bloch band structure and spatial acoustic attenuation are presented in Fig.~\ref{fig:fig1_baseline_dispersion} and Fig.~\ref{fig:fig2_baseline_attenuation}, comparing the uncoupled mechanical conservative baseline ($\beta \to 0$) and the active DPL thermoelastic metamaterial.

\begin{figure}[htbp]
\centering
\includegraphics[width=0.82\textwidth]{figures/fig1_baseline_dispersion.png}
\caption{Baseline Bloch dispersion diagram ($\Omega$ versus $k_r a / \pi$) for Family S1, comparing the conservative mechanical baseline ($\beta \to 0$, blue circles) and the active DPL thermoelastic metamaterial (red crosses). The primary closed Bragg stop band opens at $\Omega \in [0.6687, 0.7040]$.}
\label{fig:fig1_baseline_dispersion}
\end{figure}

\begin{figure}[htbp]
\centering
\includegraphics[width=0.72\textwidth]{figures/fig2_baseline_attenuation.png}
\caption{Calibrated baseline spatial acoustic attenuation $\alpha a = |k_i a|$ versus normalized frequency $\Omega$ for Family S1. The conservative baseline exhibits numerical-floor attenuation ($\alpha a \approx 1.67 \times 10^{-5}$) across pass bands, whereas active DPL thermoelasticity establishes an active dissipation baseline ($\alpha a_{\text{mean}} = 2.21 \times 10^{-4}$).}
\label{fig:fig2_baseline_attenuation}
\end{figure}

In the dispersion spectrum (Fig.~\ref{fig:fig1_baseline_dispersion}), both conservative and DPL systems exhibit nearly identical real Bloch wavenumber trajectories $k_r a / \pi$ at low frequencies ($\Omega < 0.5$). The primary Bragg band gap opens at $\Omega_L = 0.6687$ and closes at $\Omega_U = 0.7040$, producing a closed stop band of width $\Delta\Omega = 0.0354$ and gap-to-midgap ratio $\Delta\Omega / \Omega_c = 0.0515$. A secondary stop band opens at $\Omega_L = 1.3051$.

A fundamental physical distinction emerges in the spatial attenuation spectrum (Fig.~\ref{fig:fig2_baseline_attenuation}):
\begin{itemize}
    \item \textbf{Conservative System ($\beta \to 0$):} In the propagating pass bands, spatial attenuation is governed strictly by the eigensolver numerical residual floor ($\alpha a_{\text{mean}} = 1.67 \times 10^{-5}$, with minimum value $2.08 \times 10^{-9}$). Within the Bragg stop band, $\alpha a$ surges sharply to a peak value of $4.915$, representing pure geometric wave evanescence without thermal dissipation.
    \item \textbf{Active DPL Thermoelastic System:} Irreversible thermoelastic entropy generation establishes a continuous, non-zero attenuation baseline across the entire acoustic pass band ($\alpha a_{\text{mean}} = 2.21 \times 10^{-4}$, peaking at $2.77 \times 10^{-2}$ near the band edge). Furthermore, the abrupt mathematical cusp at the zone boundary is blunted into a smooth, continuous dissipation transition zone.
\end{itemize}
This quantitative comparison demonstrates the core physical distinction: in this 1D periodic continuum, Bragg stop bands arise from geometric wave reflection and destructive interference across acoustic impedance contrasts, producing reactive spatial evanescence, whereas DPL thermoelasticity induces irreversible thermodynamic dissipation across the propagating pass bands.

\subsection{Material Contrast and Bragg Scattering Evolution}
\label{sec:res_s2}

The evolution of phononic band gaps as a function of the material contrast parameter $\chi \in [0.0, 1.0]$ is illustrated in Fig.~\ref{fig:fig3_material_contrast}.

\begin{figure}[htbp]
\centering
\includegraphics[width=0.88\textwidth]{figures/fig3_material_contrast.png}
\caption{Material contrast sweep (Family S2) showing the emergence and evolution of Bragg band gaps across $\chi = 0.0$ (identical layers), $\chi = 0.5$, and $\chi = 1.0$ (full contrast).}
\label{fig:fig3_material_contrast}
\end{figure}

In the identical-layers limit ($\chi = 0.0, A = B$), the acoustic impedance ratio is unity ($Z_B / Z_A = 1.0$). Consequently, Bragg scattering is completely extinguished, and the band-gap width is identically zero:
\begin{equation}
\Delta\Omega \equiv 0.0000 \quad (\chi = 0.0).
\end{equation}
We emphasize an essential theoretical principle: \textbf{identical layers suppress material-contrast-induced Bragg gaps, but do not eliminate intrinsic gradient-elastic dispersion or thermal attenuation}. In the homogeneous medium, higher-order strain gradients still induce dispersion, and DPL thermoelasticity maintains a non-zero attenuation baseline ($\alpha a_{\text{mean}} = 3.66 \times 10^{-3}$).

As material contrast increases to $\chi = 0.5$, impedance mismatch opens an intermediate Bragg gap at $\Omega \in [0.9162, 1.0045]$ ($\Delta\Omega = 0.0884$). At full contrast ($\chi = 1.0$), strong acoustic mismatch ($Z_B / Z_A \approx 0.0935$) widens the stop band structure, shifting Gap 1 to $\Omega \in [0.6687, 0.7040]$ and opening Gap 2 at $\Omega_L = 1.3051$.

\subsection{Dipolar Gradient Length Scales: Micro-Stiffness vs.\ Micro-Inertia}
\label{sec:res_s3}

Fig.~\ref{fig:fig4_gradient_lengths} presents the parametric sensitivity of acoustic dispersion to the dipolar gradient length scales $d_1/a$ and $\sqrt{c_1}/a$.

\begin{figure}[htbp]
\centering
\includegraphics[width=0.88\textwidth]{figures/fig4_gradient_lengths.png}
\caption{Dipolar gradient-elastic length-scale sensitivity on continuous acoustic dispersion (Family S3): (a) micro-inertia variation $d_1/a \in \{0.1, 0.5, 1.0\}$; (b) micro-stiffness variation $\sqrt{c_1}/a \in \{0.1, 0.5, 0.8\}$.}
\label{fig:fig4_gradient_lengths}
\end{figure}

The two dipolar length scales exhibit diametrically opposing physical influences:
\begin{itemize}
    \item \textbf{Micro-Inertia ($d_1/a$):} Kinetic micro-inertia introduces dynamic mass penalty at higher frequencies, causing dispersive softening. Increasing $d_1/a$ from $0.1$ to $1.0$ reduces the phase velocity at short wavelengths and shifts the primary Bragg band gap downward from $\Omega \in [0.7040, 1.1990]$ to $\Omega \in [0.4035, 0.6157]$ ($\Delta\Omega = 0.2121$).
    \item \textbf{Micro-Stiffness ($\sqrt{c_1}/a$):} Higher-order strain gradients enhance deformation resistance, stiffening the acoustic response. Increasing $\sqrt{c_1}/a$ from $0.1$ to $0.8$ elevates phase velocity and shifts the primary band gap upward from $\Omega \in [0.2444, 0.6157]$ to $\Omega \in [0.8101, 1.0399]$.
    \item \textbf{Classical Elastic Limit ($c, d \to 0$):} When gradient parameters approach zero ($c_1 = 10^{-10}\text{ m}^2, d_1 = 10^{-5}\text{ m}$), higher-order dispersion vanishes, recovering classical acoustic band structure.
\end{itemize}

\subsection{Filling Fraction and Geometric Tuning}
\label{sec:res_s4}

Fig.~\ref{fig:fig5_filling_fraction} depicts the influence of geometric layer thickness asymmetry $\eta = a_1/a \in \{0.2, 0.5, 0.8\}$.

\begin{figure}[htbp]
\centering
\includegraphics[width=0.82\textwidth]{figures/fig5_filling_fraction.png}
\caption{Filling fraction sweep (Family S4) illustrating the effect of layer thickness ratio $\eta = a_1/a$ on band-gap placement and width.}
\label{fig:fig5_filling_fraction}
\end{figure}

For thin Layer A ($\eta = 0.2$), the unit cell is dominated by Layer B (Aluminum benchmark, which has a lower phase speed in this normalized system), shifting the primary band gap to $\Omega \in [0.5626, 0.7394]$ ($\Delta\Omega = 0.1768$). Conversely, for thick Layer A ($\eta = 0.8$), Epoxy dominance stiffens the effective wave speed, shifting band edges toward higher frequencies. The symmetric configuration ($\eta = 0.5$) provides optimal destructive interference balance for mid-frequency isolation.

\subsection{Dual-Phase-Lag (DPL) Thermal Time Lags}
\label{sec:res_s5}

The influence of Tzou's non-Fourier time lags on acoustic dissipation is shown in Fig.~\ref{fig:fig6_dpl_lags}.

\begin{figure}[htbp]
\centering
\includegraphics[width=0.88\textwidth]{figures/fig6_dpl_lags.png}
\caption{Dual-Phase-Lag (DPL) non-Fourier thermal time lags on continuous acoustic attenuation (Family S5): (a) heat flux relaxation lag $\tau_q \in \{1\text{ ps}, 10\text{ ps}, 1\text{ ns}\}$; (b) temperature gradient retardation lag $\tau_\theta \in \{0.1\text{ ps}, 2\text{ ps}, 100\text{ ps}\}$.}
\label{fig:fig6_dpl_lags}
\end{figure}

At ultrasonic frequencies, when the relaxation lag is small ($\tau_q = 1$ to $10\text{ ps}$), thermal transport remains in the quasi-diffusive regime, with pass-band attenuation averaging $\alpha a \approx 1.69 \times 10^{-5}$. When $\tau_q$ increases to $1\text{ ns}$, the thermal Deborah number $W = \omega \tau_q$ approaches $\mathcal{O}(10^{-2})$, activating hyperbolic thermal wave transport. The finite speed of heat propagation induces a phase lag between conductive heat flux and temperature gradient, elevating the pass-band attenuation floor to $\alpha a_{\text{mean}} = 1.70 \times 10^{-4}$ and peak attenuation to $2.19 \times 10^{-2}$. Conversely, increasing the retardation lag $\tau_\theta$ accelerates thermal conductivity relaxation, smoothing out high-frequency attenuation peaks.

\subsection{Thermoelastic Coupling Intensity and Band-Edge Blunting}
\label{sec:res_s6}

Fig.~\ref{fig:fig7_thermoelastic_coupling} demonstrates the sensitivity of dispersion and attenuation to the thermoelastic coupling coefficient $\alpha_t / \alpha_{t,\text{base}} \in \{0.0, 0.5, 1.0, 2.0\}$.

\begin{figure}[htbp]
\centering
\includegraphics[width=0.88\textwidth]{figures/fig7_thermoelastic_coupling.png}
\caption{Thermoelastic coupling intensity sensitivity and band-edge blunting (Family S6): (a) acoustic dispersion modification; (b) spatial attenuation floor scaling across $\alpha_t / \alpha_{t,\text{base}} \in \{0.0, 0.5, 1.0, 2.0\}$.}
\label{fig:fig7_thermoelastic_coupling}
\end{figure}

In the uncoupled limit ($\alpha_t = 0, \beta \to 0$), pass-band attenuation is governed by the numerical noise floor ($\sim 10^{-9}$), and dispersion curves display sharp slope discontinuities at the Brillouin zone boundary ($k_r a / \pi = 1.0$). As coupling increases ($0.5\times \to 2.0\times$), the pass-band attenuation floor increases monotonically:
\begin{equation}
\alpha a_{\text{mean}} = 1.67 \times 10^{-5} \quad (\beta \to 0) \quad \longrightarrow \quad \alpha a_{\text{mean}} = 8.84 \times 10^{-4} \quad (2.0\times \alpha_t),
\end{equation}
with peak acoustic attenuation reaching $5.53 \times 10^{-2}$. Here, $\alpha a_{\text{mean}} = 8.84 \times 10^{-4}$ represents the mean attenuation along the continuously tracked propagating acoustic branch across the investigated frequency range. By comparison, applying a strict automated pass-band filtering criterion ($\alpha a < 0.1$) yields $1.26 \times 10^{-5}$ for case \texttt{S6\_alpha20}, because the severe band-edge blunting elevates attenuation above the threshold over an extended frequency interval; these metrics reflect distinct statistical definitions rather than contradictory results. Crucially, the non-conservative thermal damping blunts the sharp mathematical cusps of Bragg band edges into smooth, continuous dissipation zones.

\subsection{Combined Parameter Interactions}
\label{sec:res_s7}

Fig.~\ref{fig:fig8_combined_interaction} displays the factorial interaction among all four governing mechanisms across six representative combinations in Family S7.

\begin{figure}[htbp]
\centering
\includegraphics[width=0.88\textwidth]{figures/fig8_combined_interaction.png}
\caption{Factorial combined parameter interaction study (Family S7) across six distinct mechanism pairings, confirming the independent physical action of Bragg scattering, gradient elasticity, and DPL thermoelasticity.}
\label{fig:fig8_combined_interaction}
\end{figure}

The factorial analysis demonstrates that the structural response cannot be modeled by simple linear superposition. While Bragg scattering dictates the frequency placement of stop bands, dipolar micro-stiffness and micro-inertia independently modulate group velocity and branch curvature, and DPL thermoelasticity superimposes a broadband dissipation envelope.

\subsection{Quantitative Synthesis and Decoupling Map}
\label{sec:res_synthesis}

Fig.~\ref{fig:fig9_bandgap_summary} summarizes the evolution of Bragg band-gap widths as a function of material contrast $\chi$ and filling fraction $\eta$, extracted from the authoritative dataset `PRODUCTION\_BANDGAP_SUMMARY.csv`.

\begin{figure}[htbp]
\centering
\includegraphics[width=0.92\textwidth]{figures/fig9_bandgap_summary.png}
\caption{Authoritative Bragg band-gap summary (Family S2 and S4). Gap 1 ($\Omega \in [0.67, 0.70]$) is a closed Bragg band gap. Gap 2 ($\Omega_L = 1.3051$) is explicitly annotated as an open band gap whose upper edge lies beyond the investigated frequency ceiling ($\Omega = 1.80$).}
\label{fig:fig9_bandgap_summary}
\end{figure}

\begin{figure}[htbp]
\centering
\includegraphics[width=0.75\textwidth]{figures/fig10_synthesis_map.png}
\caption{Scientific synthesis map (Family S1 and S3) demonstrating the tripartite decoupling among: (1) Pure Bragg scattering ($\beta \to 0$, gradient elastic); (2) Classical thermoelasticity ($c, d \to 0$); and (3) Full coupled DPL dipolar gradient metamaterial.}
\label{fig:fig10_synthesis_map}
\end{figure}

\textbf{Crucial Band-Gap Demarcation:} In the production dataset, 23 gap records terminate at the frequency ceiling $\Omega = 1.8000$ and are flagged with \texttt{is\_boundary\_truncated = True}. These records (including baseline Gap 2, $\Omega_L = 1.3051$) represent \textbf{open stop bands whose upper physical band edge lies beyond the investigated frequency range}. They must not be interpreted as closed physical band edges. Table~\ref{tab:bandgap_summary} presents the authoritative band-gap boundaries.

\begin{table}[htbp]
\centering
\caption{Extracted Bragg band-gap summary for representative production sweep cases.}
\label{tab:bandgap_summary}
\small
\begin{tabular*}{\textwidth}{@{\extracolsep{\fill}}llccccc@{}}
\toprule
Case ID & Physical Description & Gap \# & $\Omega_L$ & $\Omega_U$ & $\Delta\Omega$ & Boundary Truncated? \\
\midrule
\texttt{S1\_cons} & Conservative Baseline ($\beta \to 0$) & 1 & $1.5702$ & $1.6763$ & $0.1061$ & False (Closed) \\
\texttt{S1\_dpl} & Active DPL Baseline & 1 & $0.6687$ & $0.7040$ & $0.0354$ & False (Closed) \\
\texttt{S1\_dpl} & Active DPL Baseline & 2 & $1.3051$ & $>1.8000$ & Open & \textbf{True (Open at boundary)} \\
\texttt{S2\_chi00} & Identical Layers Limit ($\chi=0.0$) & --- & --- & --- & $\mathbf{0.0000}$ & False (No gap) \\
\texttt{S2\_chi05} & Intermediate Contrast ($\chi=0.5$) & 1 & $0.9162$ & $1.0045$ & $0.0884$ & False (Closed) \\
\texttt{S2\_chi10} & Full Contrast Baseline ($\chi=1.0$) & 1 & $0.6687$ & $0.7040$ & $0.0354$ & False (Closed) \\
\texttt{S2\_chi10} & Full Contrast Baseline ($\chi=1.0$) & 2 & $1.3051$ & $>1.8000$ & Open & \textbf{True (Open at boundary)} \\
\texttt{S3\_d01} & Low Micro-Inertia ($d_1/a=0.1$) & 1 & $0.7040$ & $1.1990$ & $0.4949$ & False (Closed) \\
\texttt{S3\_d10} & High Micro-Inertia ($d_1/a=1.0$) & 1 & $0.4035$ & $0.6157$ & $0.2121$ & False (Closed) \\
\texttt{S3\_c01} & Low Micro-Stiffness ($\sqrt{c_1}/a=0.1$) & 1 & $0.2444$ & $0.6157$ & $0.3712$ & False (Closed) \\
\texttt{S3\_c08} & High Micro-Stiffness ($\sqrt{c_1}/a=0.8$) & 1 & $0.8101$ & $1.0399$ & $0.2298$ & False (Closed) \\
\texttt{S3\_classical} & Classical Limit ($c, d \to 0$) & 1 & $0.2091$ & $0.7571$ & $0.5480$ & False (Closed) \\
\texttt{S4\_eta02} & Asymmetric Filling ($\eta=0.2$) & 1 & $0.5626$ & $0.7394$ & $0.1768$ & False (Closed) \\
\bottomrule
\end{tabular*}
\end{table}

The final synthesis map (Fig.~\ref{fig:fig10_synthesis_map}) explicitly decouples the three mechanisms:
\begin{enumerate}
    \item \textbf{Mechanism 1 (Pure Bragg Scattering):} Conservative wave reflection at periodic interfaces creates stop bands without intrinsic dissipation.
    \item \textbf{Mechanism 2 (Gradient Elastic Dispersion):} Micro-inertia and micro-stiffness alter acoustic phase and group velocities, shifting cut-offs and governing branch curvature.
    \item \textbf{Mechanism 3 (Dual-Phase-Lag Dissipation):} Non-Fourier thermal lagging introduces irreversible thermodynamic dissipation across pass bands and blunts band edges.
\end{enumerate}

---

\section{Numerical Conditioning and Algorithmic Robustness}
\label{sec:conditioning}

A rigorous assessment of numerical stability was conducted throughout the 3,600 unit-cell solves.
\begin{itemize}
    \item \textbf{Equilibrated Modal Conditioning:} After canonical two-sided equilibration ($P_{\text{equil}} = D_{\text{row}}^{-1} P D_{\text{col}}^{-1}$), the modal condition numbers remain exceptionally well-conditioned: $\kappa(P_{\text{equil}}) \le 22.72$ across all standard baseline sweeps. In the classical elastic limit case (\texttt{S3\_classical}), $\kappa(P_{\text{equil}})$ peaks at $4648.99$ at $\Omega = 0.05$ and relaxes to $129.15$ at $\Omega = 1.80$.
    \item \textbf{Raw Condition Number Context:} The unscaled condition number $\kappa(P_{\text{raw}}) \sim 10^{15} - 10^{20}$ is purely an artifact of disparate SI physical units ($10^{-9}\text{ m}$ displacement versus $10^{11}\text{ Pa}$ stress versus $10^{15}\text{ W/m}^2$ heat flux). It is not the production stability metric.
    \item \textbf{Evanescent Mode Filtering in the Classical Limit:} In case \texttt{S3\_classical}, gradient length scales are reduced to $c_1 = 10^{-10}\text{ m}^2$. Consequently, the boundary layer penetration depth $\delta \sim \sqrt{c}$ vanishes. Across layer thicknesses of $5\text{ mm}$, the corresponding boundary layer modes undergo decay exceeding $e^{-500,000}$, underflowing double-precision floating-point arithmetic ($|\lambda| < 10^{-15}$). Filtering these 223 underflow records reflects physical asymptotic boundary layer condensation, not modal disappearance.
\end{itemize}

---

\section{Limitations and Scope Constraints}
\label{sec:limitations}

To ensure rigorous scientific integrity, several inherent scope constraints and limitations of the present study are explicitly acknowledged:
\begin{enumerate}
    \item \textbf{Finite Frequency Window:} The production campaign was bounded to $\Omega \in [0.05, 1.80]$. Consequently, 23 band gaps terminating at $\Omega = 1.8000$ remain open at the upper boundary. Evaluating their upper closure frequencies requires extended high-frequency sweeps.
    \item \textbf{Evanescent Float64 Underflow:} Extreme boundary-layer modes in classical limit transitions undergo decay below double-precision machine noise ($|\lambda| < 10^{-15}$), requiring quadrupled precision or asymptotic boundary layer matching for sub-nanometer scale features.
    \item \textbf{Idealized Layer Interfaces:} Interfaces are modeled as mechanically and thermally ideal. Real layered metamaterials exhibit Kapitza thermal boundary resistance and imperfect mechanical bonding, which introduce additional acoustic impedance and thermal jumps.
    \item \textbf{Linearized Kinematics and Thermoelasticity:} The governing equations assume small-strain kinematics and linear thermoelastic constitutive laws, omitting finite deformations and temperature-dependent material properties.
    \item \textbf{One-Dimensional Architecture:} The present formulation focuses on 1D normal-incidence acoustic propagation. Oblique wave incidence and two- or three-dimensional phononic lattices introduce complex polarization coupling and directional band-gap anisotropy.
\end{enumerate}

---

\section{Conclusions}
\label{sec:conclusions}

This study developed a complete, analytically derived, and unconditionally stable transfer-matrix framework for one-dimensional periodic metamaterials governed by Mindlin Form-II dipolar gradient elasticity coupled with Tzou's Dual-Phase-Lag (DPL) heat conduction. The formulation incorporates self-consistent generalized boundary tractions derived from Hamilton's variational principle. By transforming interface matching and Bloch periodicity into a bounded Generalized Interface Eigenvalue Problem ($A\mathbf{c} = \lambda B\mathbf{c}$) with canonical two-sided equilibration ($\kappa(P_{\text{equil}}) \le 22.72$), numerical exponential overflow from stiff thermal diffusion modes was unconditionally eliminated.

The numerical implementation was validated against the independent analytical benchmark of Papargyri-Beskou et al. \cite{papargyribeskou2009wave} to machine precision ($2.45 \times 10^{-15}$ relative error) and against classical phononic crystal band gaps. A comprehensive parametric production campaign encompassing 36 cases across seven sweep families (17,747 modal records) yielded the following fundamental conclusions:
\begin{enumerate}
    \item \textbf{Bragg Stop Bands vs.\ DPL Dissipation:} Spatial acoustic attenuation within Bragg band gaps is governed by geometric wave evanescence ($\alpha a \approx 4.915$), whereas DPL thermoelasticity establishes an active thermodynamic dissipation floor ($\alpha a \sim 2.2 \times 10^{-4}$ to $5.5 \times 10^{-2}$) across all propagating pass bands.
    \item \textbf{Material Contrast Independence:} In the identical-layers limit ($\chi = 0$), acoustic impedance contrast vanishes and Bragg band gaps are completely eliminated ($\Delta\Omega \equiv 0.0000$). However, intrinsic dipolar gradient dispersion and thermal attenuation remain active throughout the homogeneous medium.
    \item \textbf{Opposing Gradient Scale Modulations:} Dipolar micro-stiffness ($\sqrt{c}/a$) induces dispersive stiffening, increasing acoustic phase velocity and shifting band-gap edges to higher frequencies. Conversely, kinetic micro-inertia ($d/a$) induces dispersive softening, lowering phase velocity and shifting band gaps downward.
    \item \textbf{Non-Fourier Thermal Lag Effects:} Increasing the heat flux relaxation lag $\tau_q$ elevates pass-band attenuation and blunts sharp Bragg band edges into smooth, continuous dissipation zones, while the temperature gradient retardation lag $\tau_\theta$ smooths high-frequency attenuation peaks.
\end{enumerate}
These findings provide quantitative design principles for multifunctional microstructured metamaterials requiring simultaneous acoustic vibration isolation, tailored dispersion, and passive thermal wave attenuation. Future investigations will extend this formulation to two-dimensional phononic lattices, interface thermal contact resistance, and nonlinear thermoelastic kinematics.

---

\section*{Replication and Open Data Statement}
All production datasets (`.csv`, `.json`), calibrated figure generation scripts, modular solver packages, and verification suites are permanently archived and directly reproducible in the repository workspace.

---

\bibliographystyle{unsrt}
\bibliography{references}

\end{document}
